% Auto-generated translated glossary artifact
\newacronym{2DEG}{2DEG}{Zweidimensionales Elektronengas}
\newacronym{3AC}{3AC}{dreiphasiger Wechselstrom}
\newglossaryentry{AC}
{
  name={AC},
  description={Wechselstrom (Alternating Current)}
}
\newacronym{AD}{AD}{Analog Digital}
\newacronym{ADC}{ADC}{Analog Digital Wandler}
\newacronym{AFE}{AFE}{Active Front End}
\newacronym{ASM}{ASM}{Asynchronmaschine}
\newacronym{AlGaN}{AlGaN}{Aluminiumgalliumnitrid}
\newacronym{AlN}{AlN}{Aluminiumnitrid}
\newacronym{BLDC}{BLDC}{Bürstenloser Gleichstrommotor}
\newglossaryentry{BNC}
{
  name={BNC},
  description={Bayonet-Neill-Concelman Bajonettverschluss}
}
\newacronym{C}{C}{Kohlenstoff}
\newacronym{DBS}{DBS}{Drehstrombrückenschaltung}
\newacronym{DC}{DC}{Gleichstrom (Direct Current)}
\newacronym{DC/DC-Wandler}{DC/DC-Wandler}{Gleichspannungswandler}
\newacronym{DG}{DG}{Diodengleichrichter}
\newacronym{DMOS}{DMOS}{double-diffused metal-oxide semiconductor}
\newglossaryentry{DSP}
{
  name={DSP},
  description={digitaler Signalprozessor}
}
\newglossaryentry{DUT}
{
  name={DUT},
  description={Device Under Test}
}
\newacronym{E-HEMT}{E-HEMT}{enhancement-mode high-electron-mobility-transistor}
\newglossaryentry{EEPROM}
{
  name={EEPROM},
  description={electrically erasable programmable read-only memory}
}
\newacronym{EMV}{EMV}{Elektromagnetischen Verträglichkeit}
\newglossaryentry{EPSR}
{
  name={EPSR},
  description={Einplatinenstromrichter}
}
\newacronym{ESB}{ESB}{Ersatzschaltbild}
\newacronym{ETI}{ETI}{Elektrotechnisches Institut}
\newacronym{ETI-Bus}{ETI-Bus}{kommunikationsbus des DSP-Systems}
\newglossaryentry{FAM}
{
  name={FAM},
  description={Feldorientierte Regelung der Asynchronmaschine}
}
\newacronym{FB}{FB}{Full Bridge - Vollbrücken}
\newacronym{FPGA}{FPGA}{Field Programmable Gate Array}
\newacronym{GIT}{GIT}{Gate Injection Transistoren}
\newacronym{GM}{GM}{Gleichstrommaschine}
\newacronym{GaN}{GaN}{Galliumnitrid}
\newglossaryentry{HB}
{
  name={HB},
  description={Halbbrücke}
}
\newacronym{HEMT}{HEMT}{high-electron-mobility-transistor}
\newacronym{HIL}{HIL}{Hardware-in-the-Loop}
\newacronym{HMK}{HMK}{Hochleistungsmodulatorkarte}
\newacronym{IC}{IC}{Integrated circuit}
\newacronym{IGBT}{IGBT}{Insulated Gate Bipolar Transistor}
\newacronym{JFET}{JFET}{junction field effect transistor}
\newacronym{KIT}{KIT}{Karlsruher Institut für Technologie}
\newglossaryentry{LED}
{
  name={LED},
  description={light emitting diode}
}
\newacronym{LWL}{LWL}{Lichtwellenleiter}
\newacronym{M3C}{M3C}{Modularer-Multilevel-Matrix-Umrichter}
\newacronym{MAC}{MAC}{Multiply-Accumulate}
\newacronym{MMC}{MMC}{Modularer-Multilevel-Umrichter}
\newglossaryentry{MOSFET}
{
  name={MOSFET},
  description={Metall Oxid Halbleiter Feldeffekttransistor}
}
\newacronym{NPT}{NPT}{Non-Punch-Through}
\newacronym{OPV}{OPV}{Operationsverstärker}
\newacronym{PCB}{PCB}{Printed circuit board}
\newacronym{PCC}{PCC}{Point of common coupling - Netzanschlusspunkt}
\newacronym{PHIL}{PHIL}{Power-Hardware-in-the-Loop}
\newacronym{PLL}{PLL}{Phase Look Loop}
\newacronym{PSM}{PSM}{Permanenterregte Synchronmaschiene}
\newacronym{PT}{PT}{Punch-Through}
\newacronym{PWM}{PWM}{Pulsweitenmodulation (Puls Wide Modulation)}
\newacronym{SIL}{SIL}{Software-in-the-Loop}
\newglossaryentry{SMD}
{
  name={SMD},
  description={Surface-mounted device}
}
\newacronym{SOA}{SOA}{Safe Operating Area}
\newacronym{Si}{Si}{Silizium}
\newacronym{SiC}{SiC}{Siliziumkarbid}
\newacronym{SoC}{SoC}{System on Chip}
\newacronym{TA}{TA}{Taktperiode des Umrichtersystems}
\newacronym{THT}{THT}{Through-hole technology}
\newglossaryentry{TSS}
{
  name={TSS},
  description={Tiefsetzsteller}
}
\newacronym{USB}{USB}{Universal Serial Bus}
\newacronym{UVLO}{UVLO}{Undervoltage-lockout}
\newacronym{VHDL}{VHDL}{Very High Speed Integrated Circuit Hardware Description Language}
\newacronym{WBG}{WBG}{wide-bandgap}
\newacronym{etc.}{etc.}{et cetera}
\newglossaryentry{sym_SiO2}
{
  name={sym:SiO2},
  description={Siliziumdioxid}
}
\newglossaryentry{sym_cce}
{
  name={sym:cce},
  description={Kollektor-Emitter-Kapazität}
}
\newglossaryentry{sym_cdc}
{
  name={sym:cdc},
  description={Zwischenkreis-Kapazität}
}
\newglossaryentry{sym_cds}
{
  name={sym:cds},
  description={Drain-Source-Kapazität}
}
\newglossaryentry{sym_cgc}
{
  name={sym:cgc},
  description={Gate-Kollektor-Kapazität}
}
\newglossaryentry{sym_cgd}
{
  name={sym:cgd},
  description={Gate-Drain-Kapazität}
}
\newglossaryentry{sym_cge}
{
  name={sym:cge},
  description={Gate-Emitter-Kapazität}
}
\newglossaryentry{sym_cgs}
{
  name={sym:cgs},
  description={Gate-Source-Kapazität}
}
\newglossaryentry{sym_ciss}
{
  name={sym:ciss},
  description={Eingangskapazität}
}
\newglossaryentry{sym_coss}
{
  name={sym:coss},
  description={Ausgangskapazität}
}
\newglossaryentry{sym_crss}
{
  name={sym:crss},
  description={Rückwirkungskapazität}
}
\newglossaryentry{sym_deltaul}
{
  name={sym:deltaul},
  description={Induzierter Spannungsabfall im Einschaltvorgang}
}
\newglossaryentry{sym_deltauls}
{
  name={sym:deltauls},
  description={Induzierte Gegespannung im Ausschaltvorgang}
}
\newglossaryentry{sym_didt}
{
  name={sym:didt},
  description={Stromgradient}
}
\newglossaryentry{sym_dudt}
{
  name={sym:dudt},
  description={Spannungsgradient}
}
\newglossaryentry{sym_eoff}
{
  name={sym:eoff},
  description={Ausschaltenergie}
}
\newglossaryentry{sym_eon}
{
  name={sym:eon},
  description={Einschaltenergie}
}
\newglossaryentry{sym_gammarrotor}
{
  name={sym:gammarrotor},
  description={Winkel des Rotorfluss der Asynchronmaschine}
}
\newglossaryentry{sym_gammarstator}
{
  name={sym:gammarstator},
  description={Winkel des Statorfluss der Asynchronmaschine}
}
\newglossaryentry{sym_ic}
{
  name={sym:ic},
  description={Kollektor-Strom}
}
\newglossaryentry{sym_ices}
{
  name={sym:ices},
  description={Kollektor Reststrom}
}
\newglossaryentry{sym_id}
{
  name={sym:id},
  description={Drainstrom}
}
\newglossaryentry{sym_idss}
{
  name={sym:idss},
  description={Drain Reststrom}
}
\newglossaryentry{sym_ig}
{
  name={sym:ig},
  description={Gatestrom}
}
\newglossaryentry{sym_ilast}
{
  name={sym:ilast},
  description={Laststrom}
}
\newglossaryentry{sym_indukhaupt}
{
  name={sym:indukhaupt},
  description={Hauptinduktivität}
}
\newglossaryentry{sym_indukrotorstreu}
{
  name={sym:indukrotorstreu},
  description={Rotorstreuinduktivität}
}
\newglossaryentry{sym_indukstatorstreu}
{
  name={sym:indukstatorstreu},
  description={Statorstreuinduktivität}
}
\newglossaryentry{sym_inneresMoment}
{
  name={sym:inneresMoment},
  description={inneres Moment der Asynchronmaschine}
}
\newglossaryentry{sym_irms}
{
  name={sym:irms},
  description={Effektivwert des Stroms}
}
\newglossaryentry{sym_irr}
{
  name={sym:irr},
  description={Diodenrückstrom}
}
\newglossaryentry{sym_irrm}
{
  name={sym:irrm},
  description={Rückstromspitze}
}
\newglossaryentry{sym_lastMoment}
{
  name={sym:lastMoment},
  description={Lastmoment der Asynchronmaschine}
}
\newglossaryentry{sym_llast}
{
  name={sym:llast},
  description={Last-Induktivität}
}
\newglossaryentry{sym_lsigma}
{
  name={sym:lsigma},
  description={parasitäre Induktivität im Kommutierungskreis}
}
\newglossaryentry{sym_magnetisierungsstrom}
{
  name={sym:magnetisierungsstrom},
  description={Magnetisierungsstrom der Asynchronmaschine}
}
\newglossaryentry{sym_mue_n}
{
  name={sym:mue_n},
  description={Beweglichkeit der Elektronen}
}
\newglossaryentry{sym_mue_p}
{
  name={sym:mue_p},
  description={Beweglichkeit der Löcher}
}
\newglossaryentry{sym_n}
{
  name={sym:n},
  description={n-dotiertes Donator Halbleitermaterial}
}
\newglossaryentry{sym_n}
{
  name={sym:n+},
  description={hoch dotiertes Donator Halbleitermaterial}
}
\newglossaryentry{sym_n}
{
  name={sym:n-},
  description={niedrig dotiertes Donator Halbleitermaterial}
}
\newglossaryentry{sym_omegamech}
{
  name={sym:omegamech},
  description={mechanische Drehfrequenz des Rotors der Asynchronmaschine}
}
\newglossaryentry{sym_ommegaeck}
{
  name={sym:ommegaeck},
  description={Eckkreisfrequenz (beginn des Feldschwächebereich)}
}
\newglossaryentry{sym_ommegarotor}
{
  name={sym:ommegarotor},
  description={Speißefrequenz des Rotorflusses der Asynchronmaschine}
}
\newglossaryentry{sym_ommegastator}
{
  name={sym:ommegastator},
  description={Speißefrequenz des Statorflusses der Asynchronmaschine}
}
\newglossaryentry{sym_ommegasyn}
{
  name={sym:ommegasyn},
  description={synchrone Kreisfrequenz}
}
\newglossaryentry{sym_p}
{
  name={sym:p},
  description={p-dotiertes Akzeptor Halbleitermaterial}
}
\newglossaryentry{sym_p}
{
  name={sym:p+},
  description={hoch dotiertes Akzeptor Halbleitermaterial}
}
\newglossaryentry{sym_p}
{
  name={sym:p-},
  description={niedrig dotiertes Akzeptor Halbleitermaterial}
}
\newglossaryentry{sym_pollparzahl}
{
  name={sym:pollparzahl},
  description={Polpaarzahl der Asynchronmaschine}
}
\newglossaryentry{sym_qrr}
{
  name={sym:qrr},
  description={Reverse recovery Speicherladung einer Diode}
}
\newglossaryentry{sym_rd}
{
  name={sym:rd},
  description={Driftwiderstand des n- Gebiets}
}
\newglossaryentry{sym_rdson}
{
  name={sym:rdson},
  description={Durchlasswiderstand}
}
\newglossaryentry{sym_reibMoment}
{
  name={sym:reibMoment},
  description={Reibmoment der Asynchronmaschine}
}
\newglossaryentry{sym_rg}
{
  name={sym:rg},
  description={Gate-Widerstand}
}
\newglossaryentry{sym_rgint}
{
  name={sym:rgint},
  description={Interner Gate-Widerstand}
}
\newglossaryentry{sym_rgoff}
{
  name={sym:rgoff},
  description={Gate-Widerstand im Ausschaltvorgang}
}
\newglossaryentry{sym_rgon}
{
  name={sym:rgon},
  description={Gate-Widerstand im Einschaltvorgang}
}
\newglossaryentry{sym_rotorflusskompl}
{
  name={sym:rotorflusskompl},
  description={komplexer Rotorfluss}
}
\newglossaryentry{sym_rotorspannungkompl}
{
  name={sym:rotorspannungkompl},
  description={komplexe Rotorspannung}
}
\newglossaryentry{sym_rotorstromkompl}
{
  name={sym:rotorstromkompl},
  description={komplexer Rotorstrom}
}
\newglossaryentry{sym_rotorwiderstand}
{
  name={sym:rotorwiderstand},
  description={Rotorwiderstand}
}
\newglossaryentry{sym_rw}
{
  name={sym:rw},
  description={Wannenwiderstand}
}
\newglossaryentry{sym_statorflusskompl}
{
  name={sym:statorflusskompl},
  description={komplexer Statorfluss}
}
\newglossaryentry{sym_statorspannungkompl}
{
  name={sym:statorspannungkompl},
  description={komplexe Statorspannung}
}
\newglossaryentry{sym_statorstromkompl}
{
  name={sym:statorstromkompl},
  description={komplexer Statorstrom}
}
\newglossaryentry{sym_statorwiderstand}
{
  name={sym:statorwiderstand},
  description={Statorwiderstand 123}
}
\newglossaryentry{sym_ta}
{
  name={sym:ta},
  description={Einschaltzeit des ersten Pulses eines Doppelpulsversuchs}
}
\newglossaryentry{sym_tau}
{
  name={sym:tau},
  description={Zeitkonstante Tau}
}
\newglossaryentry{sym_td}
{
  name={sym:td},
  description={Verriegelungszeit}
}
\newglossaryentry{sym_traegheitsmoment}
{
  name={sym:traegheitsmoment},
  description={Trägheitsmoment des mechanischen Systems}
}
\newglossaryentry{sym_uce}
{
  name={sym:uce},
  description={Kollektor-Emitter-Spannung}
}
\newglossaryentry{sym_ucesat}
{
  name={sym:ucesat},
  description={Kollektor-Emitter-Spannung in Sättigung}
}
\newglossaryentry{sym_udc}
{
  name={sym:udc},
  description={Zwischenkreisspannung}
}
\newglossaryentry{sym_uds}
{
  name={sym:uds},
  description={Drain-Source-Spannung}
}
\newglossaryentry{sym_udson}
{
  name={sym:udson},
  description={Drain-Source Spannung im eingeschalteten Zustand eines MOSFET}
}
\newglossaryentry{sym_ufnull}
{
  name={sym:ufnull},
  description={Flussspannung der parasitären MOSFET body-Diode}
}
\newglossaryentry{sym_uge}
{
  name={sym:uge},
  description={Gate-Emitter-Spannung}
}
\newglossaryentry{sym_ugeoff}
{
  name={sym:ugeoff},
  description={Gate-Emitter Spannung im ausgeschalteten Zustand}
}
\newglossaryentry{sym_ugeon}
{
  name={sym:ugeon},
  description={Gate-Emitter Spannung im eingeschalteten Zustand}
}
\newglossaryentry{sym_ugeth}
{
  name={sym:ugeth},
  description={Threshold-Spannung eines IGBT}
}
\newglossaryentry{sym_ugoff}
{
  name={sym:ugoff},
  description={Ausgangsspannungpegel eines Gatetreiber zum Ausschalten eines Transistors}
}
\newglossaryentry{sym_ugon}
{
  name={sym:ugon},
  description={Ausgangsspannungpegel eines Gatetreiber zum Einschalten eines Transistors}
}
\newglossaryentry{sym_ugs}
{
  name={sym:ugs},
  description={Gate-Source-Spannung}
}
\newglossaryentry{sym_ugsoff}
{
  name={sym:ugsoff},
  description={Gate-Source Spannung im ausgeschalteten Zustand}
}
\newglossaryentry{sym_ugson}
{
  name={sym:ugson},
  description={Gate-Source Spannung im eingeschalteten Zustand}
}
\newglossaryentry{sym_ugsth}
{
  name={sym:ugsth},
  description={Threshold-Spannung eines MOSFET}
}
\newglossaryentry{sym_ugtreiber}
{
  name={sym:ugtreiber},
  description={Steuerspannung am Ausgang eines Gatetreibers}
}
\newglossaryentry{sym_uth}
{
  name={sym:uth},
  description={Threshold-Spannung}
}
\newglossaryentry{sym_uuvlods}
{
  name={sym:uuvlods},
  description={angegebener UVLO Schwellwert im Datenblatt eines Gatetreibers}
}
\newglossaryentry{sym_uuvlogan}
{
  name={sym:uuvlogan},
  description={Effektive undervoltage-lockout Schwelle des GaN-Treibers}
}
\newacronym{usw.}{usw.}{und so weiter}
\newacronym{z.B.}{z.B.}{zum Beispiel}
